\section{Introduction}
\label{intro}

Let $G=(V,E)$ be a finite connected undirected graph with positive real
edge lengths. Its graph metric $d_G$ is the shortest-path distance on $V$.
An embedding of a finite metric space $(V,d)$ into real $L_1$ has
\emph{distortion at most $C$} if, after multiplication by a positive
scalar, its distances lie between $d$ and $Cd$.
The planar embedding conjecture asks whether one constant $C$ works for
every planar graph metric, independently of the number of vertices and
the edge lengths. We prove that it does.

\begin{theorem}\label{thm:main}
There is a finite universal constant $C\ge1$ with the following property.
For every finite connected planar graph $G=(V,E)$ with strictly positive real
edge lengths, there exist a measure space $(\Omega,\mathcal F,\mu)$ and a map
$f:V\to L_1(\Omega,\mathcal F,\mu;\R)$ such that
\[
 d_G(x,y)\le\|f(x)-f(y)\|_1\le C d_G(x,y)
 \qquad(x,y\in V).
\]
\end{theorem}

This proves the planar case of the conjecture of Gupta, Newman,
Rabinovich, and Sinclair, which predicts uniformly bounded $L_1$
distortion in every proper minor-closed family~\cite{GNRS04}.
In particular, the bound is independent of the number of distance
scales in the graph metric.

\subsection{Cuts, multicommodity flow, and the embedding problem}
\label{intro:flow}

The relationship with cuts explains why $L_1$ is a natural target.
For $S\subseteq V$, put
$\delta_S(x,y)=|\1_S(x)-\1_S(y)|$.
A nonnegative combination
\begin{equation}\label{intro:cut-metric}
 D(x,y)=\sum_{S\subseteq V}w_S\delta_S(x,y),\qquad w_S\ge0,
\end{equation}
is the $\ell_1$ distance between the vectors
$(w_S\1_S(x))_{S\subseteq V}$ and $(w_S\1_S(y))_{S\subseteq V}$.
Conversely, threshold sweeps through real coordinates give such a
representation for every $L_1$ metric on a finite set. Thus we seek
weighted cuts that separate each pair in proportion to its graph
distance while charging each edge at most a constant times its length.
The edge bound then extends to every pair along a shortest path.

Linial, London, and Rabinovich established the connection between
metric embeddings and multicommodity flow~\cite{LLR95}; Aumann and
Rabani independently obtained a general approximate max-flow/min-cut
theorem~\cite{AR98}. To state the corresponding planar consequence,
give the edges nonnegative capacities and the unordered pairs of
distinct vertices nonnegative demands, not all zero. Let $\lambda_*$
be the largest common multiplier of the demands that can be routed
fractionally. Let $\phi_*$ be the minimum ratio of crossing capacity
to separated demand over cuts with positive separated demand. For
$\lambda_*>0$, the ratio $\phi_*/\lambda_*$ is the \emph{flow--cut gap}.
Gupta et al. showed that its worst value on a fixed graph equals the
worst $L_1$ distortion over all weighted shortest-path metrics on that
graph~\cite[Section~3, Theorem~3.2]{GNRS04}. Hence
Theorem~\ref{thm:main} also gives a universal bound for the undirected
fractional flow--cut gap on planar graphs, with arbitrary demands.

\subsection{Earlier results and the remaining obstacle}

The general preceding bound for $n$-vertex planar metrics is
$O(\sqrt{\log n})$, due to Rao~\cite{Rao99}. His construction embeds
first into Euclidean space and then into $L_1$; it builds on the
decomposition methods of Klein, Plotkin, and Rao for graphs excluding
a fixed minor~\cite{KPR93}. Newman and Rabinovich showed that even
series-parallel metrics can require $\Omega(\sqrt{\log n})$ distortion
into Euclidean space~\cite{NR03}. Thus a constant planar $L_1$ bound
requires a different route. Abraham, Filtser, Gupta, and Neiman later
recovered Rao's planar bound through recursive shortest-path
decompositions~\cite{AFGN22}.

Constant distortion was known under several topological restrictions.
Gupta et al. proved it for series-parallel graphs~\cite{GNRS04}.
Chakrabarti, Jaffe, Lee, and Vincent obtained the sharp upper bound
$2$~\cite{CJLV08}, with a matching lower bound due to Lee and
Raghavendra~\cite{LR10}. Chekuri, Gupta, Newman, Rabinovich, and
Sinclair gave a bound exponential in the outerplanarity
parameter~\cite{CGNRS06}: the number of rounds needed to delete all
vertices by successive outer-face removals in a planar drawing.
This gives a uniform bound when that parameter is fixed.

Another line of work restricts the pairs whose distances must be
preserved. The theorem of Okamura and Seymour, and its metric
formulation by Hurkens, Schrijver, and Tardos, give an isometric
$L_1$ representation of the metric on vertices of one
face~\cite{OS81,HST88}. If a terminal set is covered by $\gamma$
faces in a fixed drawing, Krauthgamer, Lee, and Rika obtained a
logarithmic bound in $\gamma$~\cite{KLR19}. Filtser improved this to
$O(\sqrt{\log(\gamma+1)})$~\cite{Filtser25}. Kumar constructed a
single nonexpansive map on all vertices that preserves every
cofacial pair within a universal factor, even when different pairs
lie on different faces~\cite{Kumar25}.

Geometric restrictions lead to a different family of results.
Sidiropoulos proved constant distortion for planar metrics realized
in simply connected surfaces of nonpositive curvature in the sense
of Busemann~\cite{Sidiropoulos13}. Chalopin, Chepoi, and Naves proved
that Busemann surfaces themselves embed isometrically into
$L_1$~\cite{CCN15}. These are geodesic spaces homeomorphic to the
plane whose distance is convex along pairs of geodesics.

The difficulty is to combine separation at different scales while
keeping a uniform upper bound. Sidiropoulos addressed this difficulty
by evolving a distribution of monotone cuts from coarse scales to
fine scales~\cite[Sections~4--5]{Sidiropoulos13}. Our construction
also preserves monotonicity through successive local changes.
For unrestricted planar metrics, two additional mechanisms provide
the needed control: component separation limits which pairs need
further updates, and interpolation makes local formulas agree
exactly at their interfaces. The first yields a packing bound only
for the surviving pairs; the second lets those local updates be
combined without a loss proportional to the number of scales.

The structural reductions of Lee and Sidiropoulos relate the planar
problem to the larger GNRS conjecture and a separate problem of
closure under graph sums~\cite{LS09}. Together with the
bounded-treewidth companion~\cite[Theorem~1.1]{Treewidth108}, our
theorem also gives constant distortion on their fixed-parameter
almost-embeddable families; see Corollary~\ref{cor:almost-embeddable}.
The general clique-sum closure required for the full GNRS conjecture
is not established here.

\subsection{The construction and its main ideas}

We construct two families of measurable label sets, denoted by
$\mathcal A_v$ and $\mathcal B_v$ at a point $v$.
A label determines the cut of points whose sets contain it, so
$\mu(\mathcal A_x\mathbin\triangle\mathcal A_y)$ is the resulting
cut distance. The random construction chooses these entire families;
the labels index their cuts. Every outcome has the required upper
bound. Random local changes supply the lower bound.

\paragraph{Columns and the upper bound.}
Section~\ref{geo:section} represents the geometric graph isometrically
by vertical columns with zero-cost switches at specified heights.
Choose a shortest-path tree rooted at a vertex and use distance from
the root as height. Replace each non-tree edge by two ascending
branches of the same total length and identify their tips. The
plane embedding gives a circular order of columns in which the
switching links do not cross. Figure~\ref{fig:column-realization}
illustrates this realization.

For a fixed label, the active points on each column form an upper
interval. Switches preserve membership, and the measure of labels
whose membership changes along a vertical segment equals its length.
Consequently every path bounds the cut distance between its endpoints,
regardless of the number of updates. Section~\ref{sets:section}
formulates this invariant for both systems. The problem is to obtain
lower separation while retaining it.

\paragraph{Common local formulas.}
Local updates use \emph{charts}: regions on which one scalar field
$P$ and common label variables describe the thresholds. Their form is
$h(v)+\sigma\kappa P(v)$, where $h$ is height, $\sigma$ is a fair
sign, and $\kappa>0$ is fixed. Averaging over the two signs detects
both height differences and field differences. A sufficiently small
vertical Lipschitz constant for $\kappa P$ preserves the upper-interval
property.

Within a thin height band, use distances from a basepoint to the
column midpoints to form successive radial layers. Noncrossing links
restrict the entrances to a connected region beyond a layer boundary
to neighborhoods of at most two points. This two-portal property, proved in
Section~\ref{geo:section}, underlies the interpolation in
Section~\ref{int:section}. A field with Lipschitz constant strictly
below one is approximated by pieces of the form $c\pm d_I(k,\cdot)$,
where $d_I$ is distance within the band and $k$ is one point in it.
At each entrance, the neighboring regions select exactly the same
point-distance formula, including its source and its name.
Strictness of the input slope bounds
the number of candidate formulas that can govern a small ball.
Fresh independent offsets then make a change of selected
formula across that ball unlikely.

\paragraph{The first system selects the remaining pairs.}
At scale $r$, system $A$ tries to raise the field near one endpoint
of a pair at distance comparable to $r$, leaving the other endpoint
unchanged. Section~\ref{main:section} shows that failure to produce
a definite rise supplies a short path on which the old field
decreases at almost the greatest rate its Lipschitz bound permits.

System $A$ has a second separation mechanism. For each label, refine
its active set into connected components of active columns. This
refinement retains the upper bound. A sufficiently low inactive tip
of a switching link separates active points on opposite sides and
therefore certifies a lower bound. Only pairs with descending paths
at both ends and without this certificate contribute \emph{seeds},
the vertices near which the second system may update.

Proposition~\ref{pack:bound} places the seeds in a small narrow-band
ball close to two short paths. Their distance from cell boundaries
then separates their groups along those paths, bounding the number
of groups that can affect one local update.
Figure~\ref{fig:surviving-seeds} shows the decisive case: when a path
skips a seed's column, failure of the low-tip certificate forces that
seed close to the skipping link. The ambient metric need not have
a uniform packing bound.

\paragraph{The second system separates the surviving pairs.}
The circular order supplies fields that increase along both arcs
between two column indices. Section~\ref{ord:section} constructs
such a field below a point-distance function and equal to it at a
chosen target. Equal values
can occur on opposite arcs; a signed transverse field distinguishes
the arcs and vanishes where a switch can join them.

Section~\ref{sec:secondary} uses these fields in local changes near
the seeds. Either a change can be admitted while agreeing with the
old field at the boundary, or its failure in both orientations
supplies short paths along which the ordered field rises and falls.
If this happens at both endpoints, Lemma~\ref{ord:two-side} puts them on opposite
arcs, where a transverse increment separates them. The packing
bound limits the number of competing randomized proposals, both
for the separation probability and for chart persistence.

\paragraph{Sparse schedules preserve separation.}
Section~\ref{pad:assembly} combines three estimates uniform in the
graph and the update history. At scale $r$, a point's label set moves
by at most $Mr$; an eligible pair at distance in $[r,2r)$ gains a
separation certificate with fixed positive probability; and a ball of
radius $\rho$ already in one chart in
each system is split between new charts with probability
$O(\rho/r)$. Here eligibility requires neighborhoods of both
endpoints, of radius a fixed multiple of $r$, to share the same
chart data in each system.

Within a schedule, successive scales differ by a sufficiently large
constant factor $Q$. The coarser chart-loss probabilities have a
small geometric sum, so a pair reaches its own scale with positive
probability. Later movement has another small geometric sum and
cannot erase all the gained separation. For a component certificate,
it suffices to preserve activity at the two endpoints and inactivity
at one fixed tip. Averaging finitely many interleaved schedules
covers every pair distance and completes the embedding.
